\begin{theorem}[Holonomies of a one-step flux movement]%
    \label{thm:peps_torus_two_plaquette_flux_holonomy}
    \lean{TNLean.PEPS.torusTwoPlaquetteFluxAssignment,
        TNLean.PEPS.torusTwoPlaquetteFluxAssignment_holonomy,
        TNLean.PEPS.torusTwoPlaquetteFluxAssignment_eq_treeCycleAssignment,
        TNLean.PEPS.regularWalkHolonomy_twoPlaquetteTreeCycleAssignment}
    \leanok
    \uses{thm:peps_two_plaquette_torus_geometry,
        thm:peps_regular_two_cycle_flux_move,
        def:peps_torus_plaquette_flux_geometry}
    Consider a finite torus of width at least four and height at least three.
    In its six-site rectangle $\{0,1,2\}\times\{0,1\}$, let $e_0,e_1$
    be the vertical bonds at horizontal coordinates zero and one, respectively.
    Both bonds are ordered upward. For $g\in G$, let $u_0$ assign $g$ to
    $e_0$, and let $u_1$ assign $g$ to both $e_0,e_1$; all other bonds carry
    the identity. For the counterclockwise plaquette walks based at $(0,0)$
    and $(1,0)$, their actual holonomies satisfy
    \begin{align}
        (H_L(u_0),H_R(u_0))&=(g^{-1},1), &
        (H_L(u_1),H_R(u_1))&=(1,g^{-1}).
    \end{align}
    These literal bond assignments equal the tree-cycle assignments used by
    the fixed two-cycle unitary. Thus the operation removes the left
    plaquette flux and places the same flux on the right plaquette.
    This is the endpoint calculation for the finite-torus instance
    of \cite[Theorem~6.16]{Schuch2010PEPS}; the unrestricted lattice scope
    remains separate, as recorded in \cite{gap:rmp_peps_quantum_double_g_isometry}.
\end{theorem}

\begin{proof}\leanok
    Each ordered vertical bond is traversed upward with transport $g$ and
    downward with transport $g^{-1}$. The left plaquette traverses the middle
    bond upward and the left bond downward. Each successive transport
    multiplies on the left, so its extended holonomy is $g^{-1}g=1$.
    The right plaquette traverses the middle bond downward, so its extended
    holonomy is $g^{-1}$. The initial values follow by retaining only the
    left insertion. The other sides carry the identity.
    The ordered endpoints of the two non-tree bonds are exactly the two
    specified upward vertical bonds; the tree-cycle assignment is the identity
    on every remaining bond. Hence these evaluations apply to the actual
    assignments in the physical operation.
\end{proof}
